
\chapter{DC3S: The Battery Adapter}
\label{ch:dc3s-adapter}
\label{ch:dc3s}

The Detect--Calibrate--Constrain--Shield--Certify Safety System (DC3S) is the
runtime adapter that sits at the Optimize--Dispatch boundary
(Chapter~\ref{ch:orius-context}) and translates degraded observations
into safety-guaranteed actions.  This chapter specifies the runtime objects,
derives each stage of the pipeline, and narrates a complete control step from
telemetry ingestion to certificate issuance.

% ----------------------------------------------------------------
\section{Why a Runtime Shield Is Needed}

Chapter~\ref{ch:forecasting} showed that CQR intervals under-cover
conditionally: wind coverage drops to 39--49\% in the low-reliability group,
and load coverage falls 8--20 percentage points below nominal.  A dispatch
controller that trusts these intervals without adjustment will issue actions
that violate SOC constraints on the true (unobserved) trajectory at rates far
exceeding the design budget~$\alpha$.

A runtime shield is needed because:
\begin{enumerate}
  \item Forecast calibration is a \emph{batch} property; telemetry degradation
    is a \emph{streaming} property.  No static calibration fold can anticipate
    the specific fault pattern at step~$t$.
  \item The observation quality at step~$t$ is \emph{known} to the controller
    (via packet-loss counters, staleness timers, and drift detectors) but is
    \emph{not used} by standard conformal methods.
  \item The SOC dynamics are \emph{coupled}: an unsafe action at $t$ may
    render all future actions infeasible.  The shield must prevent the
    first violation, not merely reduce its frequency.
\end{enumerate}

% ----------------------------------------------------------------
\section{Runtime Objects}
\label{sec:runtime-objects}

Table~\ref{tab:runtime-objects} defines the variables maintained by DC3S at
each control step.

\begin{table}[ht]
\centering
\caption{DC3S runtime objects at step~$t$.}
\label{tab:runtime-objects}
\begin{tabular}{cll}
\toprule
\textbf{Symbol} & \textbf{Name} & \textbf{Description} \\
\midrule
$z_t$            & True state          & Actual plant state (load, renew., SOC);
                                         unobserved at runtime \\
$\tilde{z}_t$    & Observed state      & Telemetry-degraded plant state seen by
                                         controller \\
$o_t$            & Observation vector  & Raw sensor packet (may contain NaN,
                                         delays, spikes) \\
$x_t$            & Feature vector      & Engineered features derived from
                                         $\tilde{z}_{1:t}$ \\
$w_t$            & Reliability score   & $\in [0,1]$; OQE output quantifying
                                         observation quality \\
$d_t$            & Drift flag          & $\in \{0,1\}$; Page--Hinkley detector
                                         output \\
$s_t$            & Stale count         & Non-negative integer; consecutive steps
                                         since last fresh reading \\
$C_t(\alpha)$    & Base interval       & CQR prediction interval at
                                         miscoverage~$\alpha$ \\
$C_t^{\mathrm{RAC}}(\alpha)$ & Inflated interval &
                   Reliability-adaptive conformal interval \\
$\mathcal{A}_t$  & Safe-action set     & Feasible actions under inflated
                                         uncertainty \\
$a_t^\star$      & Candidate action    & Optimiser output \\
$a_t^{\mathrm{safe}}$ & Safe action    & Projected action $\in \mathcal{A}_t$ \\
$\mathcal{C}_t$  & Certificate         & Step-level audit record \\
\bottomrule
\end{tabular}
\end{table}

% ----------------------------------------------------------------
\section{Detect: Telemetry Reliability Scoring (OQE)}
\label{sec:detect}

The Observation Quality Estimator (OQE) computes the reliability score
$w_t \in [0, 1]$ from three sub-scores:

\begin{equation}
\label{eq:oqe}
  w_t = w_t^{\mathrm{fresh}} \cdot w_t^{\mathrm{range}} \cdot w_t^{\mathrm{consist}},
\end{equation}
where:
\begin{description}
  \item[$w_t^{\mathrm{fresh}}$] Freshness score: $\exp(-\lambda_f \cdot s_t)$,
    decaying exponentially with the stale count~$s_t$.
  \item[$w_t^{\mathrm{range}}$] Range plausibility: 1 if all channels are
    within physical bounds, decayed proportionally to the number of
    out-of-range channels.
  \item[$w_t^{\mathrm{consist}}$] Temporal consistency: based on the
    Mahalanobis distance of the current reading from a short rolling window,
    penalising spikes and discontinuities.
\end{description}

The OQE computation is stateless given the rolling buffer and runs in
$<0.03$\,ms median latency
(Table~\ref{tab:dc3s-latency}, \texttt{reports/publication/dc3s\_latency\_summary.csv}).

% ----------------------------------------------------------------
\section{Calibrate: Reliability-Adaptive Interval Inflation}
\label{sec:calibrate}

The base CQR interval $C_t(\alpha)$ is inflated to compensate for the
conditional under-coverage revealed in Chapter~\ref{ch:forecasting}.  The
\emph{RAC-Cert} (Reliability-Adaptive Conformal Certificate) formula is:

\begin{equation}
\label{eq:rac-cert}
\boxed{
  C_t^{\mathrm{RAC}}(\alpha) \;=\;
  C_t(\alpha) \;\cdot\;
  \bigl[\,1 + \kappa_r\,(1 - w_t)
             + \kappa_d\,d_t
             + \kappa_s\,s_t\,\bigr]
}
\end{equation}

where:
\begin{description}
  \item[$\kappa_r$] Reliability sensitivity: scales the inflation with the
    complement of the reliability score.  Default $\kappa_r = 0.5$.
  \item[$\kappa_d$] Drift sensitivity: additive inflation when the
    Page--Hinkley detector fires.  Default $\kappa_d = 0.3$.
  \item[$\kappa_s$] Staleness sensitivity: linear inflation per stale step.
    Default $\kappa_s = 0.05$.
\end{description}

The multiplicative form ensures that under nominal conditions
($w_t = 1$, $d_t = 0$, $s_t = 0$) the inflated interval collapses to the
base interval, imposing zero additional conservatism.

\paragraph{Inflation diagnostics.}
The locked evidence in \texttt{reports/publication/dc3s\_main\_table.csv}
reports the mean and P95 of the RAC inflation factor across all scenarios
and seeds.  Under the \texttt{drift\_combo} scenario the mean inflation
is approximately $1.02$, rising to $1.05$ at the P95 level, confirming that
the adaptive layer adds modest width only when observation quality genuinely
degrades.

% ----------------------------------------------------------------
\section{Constrain: Safe-Action Sets Under Widened Uncertainty}
\label{sec:constrain}

Given the inflated interval $C_t^{\mathrm{RAC}}(\alpha)$ with half-width
$q_t^{\mathrm{RAC}}$, the safe-action set is constructed by tightening
the SOC bounds (cf.\ Equation~\eqref{eq:tightened-bounds},
Chapter~\ref{ch:battery-dynamics}):

\begin{equation}
\label{eq:safe-action-set}
  \mathcal{A}_t = \Bigl\{\,(P^{+}, P^{-}) \;\Big|\;
    E_{\min} + q_t^{\mathrm{RAC}} \;\le\; \hat{E}_{t+1}(P^{+}, P^{-})
    \;\le\; E_{\max} - q_t^{\mathrm{RAC}},\;
    \text{power and ramp constraints hold}\,\Bigr\},
\end{equation}
where $\hat{E}_{t+1}(P^{+}, P^{-})$ is the one-step-ahead SOC estimate
under the candidate action and the observed initial SOC~$\tilde{E}_t$.

If $q_t^{\mathrm{RAC}}$ is large enough that $\mathcal{A}_t = \emptyset$,
the controller reverts to a safe-standby action
$(P^{+} = 0, P^{-} = 0)$, which is always feasible.

% ----------------------------------------------------------------
\section{Shield: Projection and Repair}
\label{sec:shield}

If the candidate action $a_t^\star$ lies outside $\mathcal{A}_t$, the shield
projects it to the nearest safe action:

\begin{equation}
\label{eq:repair}
\boxed{
  a_t^{\mathrm{safe}} = \argmin_{a \in \mathcal{A}_t}
    \|a - a_t^\star\|_2
}
\end{equation}

Because $\mathcal{A}_t$ is a convex polytope (intersection of linear
half-spaces from power bounds, ramp limits, and tightened SOC bounds), the
projection is a quadratic programme solvable in closed form for the
single-step case or via an efficient QP solver for the multi-step horizon.

The \emph{intervention rate} (IR) is the fraction of steps where
$a_t^{\mathrm{safe}} \neq a_t^\star$:
\[
  \mathrm{IR} = \frac{1}{T} \sum_{t=1}^{T}
    \mathbf{1}\!\bigl[a_t^{\mathrm{safe}} \neq a_t^\star\bigr].
\]

The locked benchmark results show that DC3S achieves IR~$\approx 2{-}8\%$
under degraded scenarios while maintaining TSVR~$= 0\%$, confirming that the
economic cost of safety is modest.

% ----------------------------------------------------------------
\section{FTIT State Tube}
\label{sec:ftit}

The Fault-Tolerant Invariant Tube (FTIT) provides a secondary safety layer
that does not depend on interval calibration.  FTIT propagates worst-case
SOC bounds forward over a $k$-step horizon:

\begin{equation}
\label{eq:ftit}
\begin{aligned}
  \overline{E}_{t+k} &= \hat{E}_t + \delta_{\max}
    + k\,\eta^{+}\,P_{\max}^{+}\,\Delta t, \\
  \underline{E}_{t+k} &= \hat{E}_t - \delta_{\max}
    - k\,P_{\max}^{-}/\eta^{-}\,\Delta t,
\end{aligned}
\end{equation}
where $\delta_{\max}$ is the maximum SOC estimation error consistent with
the current reliability score.  Any action that would cause either bound to
exit $[E_{\min}, E_{\max}]$ within the $k$-step tube is rejected.

The FTIT acts as a hard backstop: even if the RAC-Cert inflation is
insufficient (e.g.\ due to a novel fault type not captured by the
sensitivity parameters), the tube prevents multi-step SOC excursions.

% ----------------------------------------------------------------
\section{Certify: Step-Level Audit Records}
\label{sec:certify}

At the conclusion of each control step, DC3S emits a certificate
$\mathcal{C}_t$ containing:

\begin{table}[ht]
\centering
\caption{DC3S certificate fields.}
\label{tab:certificate-fields}
\begin{tabular}{ll}
\toprule
\textbf{Field} & \textbf{Content} \\
\midrule
\texttt{step}                    & Time-step index \\
\texttt{timestamp}               & UTC timestamp \\
\texttt{reliability\_w}          & Observation quality score $w_t$ \\
\texttt{drift\_flag}             & Page--Hinkley drift indicator $d_t$ \\
\texttt{inflation}               & RAC-Cert multiplier \\
\texttt{proposed\_charge\_mw}    & Candidate charge power $P_t^{+\star}$ \\
\texttt{proposed\_discharge\_mw} & Candidate discharge power $P_t^{-\star}$ \\
\texttt{safe\_charge\_mw}        & Repaired charge power \\
\texttt{safe\_discharge\_mw}     & Repaired discharge power \\
\texttt{intervened}              & Boolean: was the action repaired? \\
\texttt{intervention\_reason}    & Reason string (SOC, ramp, BMS) \\
\texttt{guarantee\_passed}       & Boolean: all invariants satisfied? \\
\texttt{soc\_after\_mwh}         & Estimated SOC after action \\
\bottomrule
\end{tabular}
\end{table}

A trace of certificate records for a representative episode is stored in
\texttt{reports/publication/dc3s\_step\_trace.csv}.  The certificate
stream is the primary input to the monitoring stage
(Chapter~\ref{ch:orius-context}) and to the half-life analysis
(Chapter~\ref{ch:cert-half-life-blackout}).

% ----------------------------------------------------------------
\section{DC3S Latency Profile}

Real-time dispatch requires that the DC3S pipeline complete within the
control cycle (typically 1\,s for sub-hourly, 60\,s for hourly).
Table~\ref{tab:dc3s-latency} reports the component-level latency from the
locked benchmark.

\begin{table}[ht]
\centering
\caption{DC3S component latency (ms).  Source:
\texttt{reports/publication/dc3s\_latency\_summary.csv}.}
\label{tab:dc3s-latency}
\begin{tabular}{lrrrrr}
\toprule
\textbf{Component} & \textbf{Mean} & \textbf{Median} & \textbf{P95} &
  \textbf{P99} & \textbf{Max} \\
\midrule
Reliability scoring          & 0.026 & 0.021 & 0.043 & 0.130 & 0.712 \\
Drift update (Page--Hinkley) & 0.000 & 0.000 & 0.001 & 0.001 & 0.160 \\
Uncertainty set build        & 0.012 & 0.011 & 0.011 & 0.047 & 0.724 \\
Action repair (projection)   & 0.002 & 0.002 & 0.002 & 0.003 & 0.210 \\
\midrule
\textbf{Full DC3S step}      & 0.042 & 0.037 & 0.041 & 0.193 & 2.416 \\
\bottomrule
\end{tabular}
\end{table}

The full DC3S step completes in under 0.2\,ms at P99, well within the
budget for any practical control cycle.  The maximum of 2.4\,ms occurs
during JIT warm-up on the first step and is not representative of
steady-state performance.

% ----------------------------------------------------------------
\section{Complete Runtime Narrative}
\label{sec:runtime-narrative}

A single DC3S control step proceeds as follows:

\begin{enumerate}
  \item \textbf{Receive observation} $o_t$ from the telemetry gateway.
    If the packet is missing (dropout), carry forward the last valid reading
    and increment~$s_t$.

  \item \textbf{Detect.}  Compute $w_t = w_t^{\mathrm{fresh}} \cdot
    w_t^{\mathrm{range}} \cdot w_t^{\mathrm{consist}}$ via OQE.
    Run the Page--Hinkley test on the residual stream to set~$d_t$.

  \item \textbf{Calibrate.}  Inflate the base CQR interval:
    $C_t^{\mathrm{RAC}}(\alpha) = C_t(\alpha) \cdot
    [1 + \kappa_r(1-w_t) + \kappa_d\,d_t + \kappa_s\,s_t]$.

  \item \textbf{Constrain.}  Construct the safe-action set $\mathcal{A}_t$
    by tightening SOC bounds with the inflated half-width.  Verify the
    FTIT tube; if the tube is violated, further restrict $\mathcal{A}_t$.

  \item \textbf{Shield.}  If $a_t^\star \notin \mathcal{A}_t$, project:
    $a_t^{\mathrm{safe}} = \argmin_{a \in \mathcal{A}_t} \|a - a_t^\star\|_2$.
    Otherwise $a_t^{\mathrm{safe}} = a_t^\star$.

  \item \textbf{Certify.}  Emit the certificate $\mathcal{C}_t$ recording
    all runtime variables, the intervention decision, and the
    guarantee-check outcome.  Forward $a_t^{\mathrm{safe}}$ to the BMS.
\end{enumerate}

% ----------------------------------------------------------------
\section{Chapter Summary}

This chapter specified the DC3S battery adapter: a five-stage runtime shield
(Detect, Calibrate, Constrain, Shield, Certify) that intercepts dispatch
actions at the Optimize--Dispatch boundary.  The OQE reliability
scorer~\eqref{eq:oqe} quantifies observation quality; the RAC-Cert
formula~\eqref{eq:rac-cert} inflates intervals adaptively; the safe-action
set~\eqref{eq:safe-action-set} tightens SOC bounds accordingly; and the
projection repair~\eqref{eq:repair} moves unsafe actions to the nearest
feasible point.  The FTIT tube provides a calibration-free safety backstop.
The full pipeline completes in under 0.2\,ms at P99, and step-level
certificates create a complete audit trail.  Empirical validation is
presented in Chapter~\ref{ch:cpsbench}.
